\documentclass[aspectratio=169,11pt]{beamer}

\usetheme{Madrid}
\usecolortheme{default}
\usefonttheme{professionalfonts}
\setbeamertemplate{navigation symbols}{}
\setbeamertemplate{footline}[frame number]

\usepackage{amsmath, amssymb, booktabs, graphicx}

\newcommand{\E}{\mathbb{E}}
\newcommand{\1}{\mathbf{1}}
\newcommand{\Var}{\operatorname{Var}}
\newcommand{\Avar}{\operatorname{Avar}}
\newcommand{\argmax}{\operatorname*{arg\,max}}
\newcommand{\argmin}{\operatorname*{arg\,min}}

\title{Structural Estimation in Practice}
\subtitle{Identification, Moments, Likelihood, Simulation, and Fit \\ Empirical Methods --- Week 7}
\author{Teaching deck draft}
\date{}

\begin{document}

\begin{frame}
  \titlepage
\end{frame}

\begin{frame}{Week 7 question}
\begin{itemize}
    \item How do researchers take a structural model to data credibly?
    \item Which variation identifies the latent primitives?
    \item Which likelihood pieces, moments, or auxiliary statistics carry the identifying power?
    \item How should fit, counterfactual discipline, and uncertainty be reported?
\end{itemize}

\medskip
Structural credibility is the chain: identification $\rightarrow$ estimation $\rightarrow$ fit $\rightarrow$ counterfactual discipline $\rightarrow$ inference.
\end{frame}

\begin{frame}{Identification in structural models}
\small
\begin{tabular}{p{0.24\linewidth} p{0.29\linewidth} p{0.35\linewidth}}
\toprule
Question & Applied object & Credibility check \\
\midrule
What is latent? & preferences, costs, beliefs, technology, risk & does the paper name the primitive? \\
What is observed? & choices, states, outcomes, prices, rules & are the data rich enough? \\
What moves it? & policy variation, prices, histories, transitions & which variation shifts the primitive? \\
What else could explain it? & omitted state, heterogeneity, expectations & what breaks observational equivalence? \\
What uses it? & welfare or policy counterfactual & is the counterfactual near support? \\
\bottomrule
\end{tabular}

\[
\widehat h_N(Y,X,S,P)
\quad \Longleftrightarrow \quad
h(\theta;\mathcal M)
\]
\end{frame}

\begin{frame}{Likelihood vs moments vs simulation}
\small
\textbf{Likelihood}
\[
\widehat\theta_{\mathrm{ML}}
=
\argmax_\theta
\sum_i \log f_\theta(y_i\mid x_i,s_i,p_i)
\]

\textbf{GMM / method of moments}
\[
\widehat\theta_{\mathrm{GMM}}
=
\argmin_\theta
\widehat g_N(\theta)'W_N\widehat g_N(\theta),
\qquad
\widehat g_N(\theta)=\widehat m^{data}-m^{model}(\theta)
\]

\textbf{SMM}
\[
\widehat\theta_{\mathrm{SMM}}
=
\argmin_\theta
\left[\widehat m^{data}-\widehat m^{sim}(\theta;R)\right]'
W_N
\left[\widehat m^{data}-\widehat m^{sim}(\theta;R)\right]
\]
\end{frame}

\begin{frame}{Indirect inference intuition}
\[
\widehat\theta_{\mathrm{II}}
=
\argmin_\theta
\left[
\widehat a^{data}
-
\widehat a^{sim}(\theta)
\right]'
W_N
\left[
\widehat a^{data}
-
\widehat a^{sim}(\theta)
\right]
\]

\begin{itemize}
    \item The structural model generates simulated histories.
    \item The auxiliary model summarizes behavior in a familiar empirical language.
    \item The auxiliary statistics must capture the behavioral margins that matter.
    \item Good auxiliary fit is not enough if the statistic is unrelated to the counterfactual.
\end{itemize}
\end{frame}

\begin{frame}{Weighting matrices and targeted moments}
\small
\[
Q_N(\theta)=\widehat g_N(\theta)' W_N \widehat g_N(\theta),
\qquad
W^\ast=S^{-1},
\qquad
S=\E[g_i(\theta_0)g_i(\theta_0)']
\]

\[
J_N
=
N\widehat g_N(\widehat\theta)'
\widehat W_N
\widehat g_N(\widehat\theta)
\]

\begin{itemize}
    \item Weighting determines which discrepancies the criterion punishes.
    \item Targeted moments define the model's empirical obligations.
    \item Overidentification lets the model fail, but non-rejection is not validation.
    \item Report moment scaling, weights, dominated moments, and sensitivity.
\end{itemize}
\end{frame}

\begin{frame}{Fit, validation, counterfactual discipline}
\small
\begin{tabular}{p{0.26\linewidth} p{0.30\linewidth} p{0.33\linewidth}}
\toprule
Object & What to show & Why it matters \\
\midrule
Targeted fit & empirical vs model moments & criterion transparency \\
Untargeted fit & distributions, transitions, holdout moments & validation beyond construction \\
Comparative statics & response to prices, policies, states & economic mechanism \\
Counterfactual distance & support, policy range, state distribution & extrapolation risk \\
Sensitivity & weights, calibrated values, seeds, starts & numerical and model dependence \\
\bottomrule
\end{tabular}

\[
\widehat\theta
\Longrightarrow
\left\{
\sigma_{\widehat\theta}(a\mid s;p),
F_{\widehat\theta}(s'\mid s,a;p),
Y_{\widehat\theta}(p),
W_{\widehat\theta}(p)
\right\}
\]
\end{frame}

\begin{frame}{Variance, bootstrap, delta method}
\small
\textbf{Likelihood / sandwich}
\[
\widehat{\Var}(\widehat\theta)
\approx
\frac{1}{N}\widehat A^{-1}\widehat B\widehat A^{-1},
\qquad
\widehat B=\frac{1}{N}\sum_i s_i(\widehat\theta)s_i(\widehat\theta)'
\]

\textbf{GMM}
\[
\Avar(\widehat\theta)
=
\frac{1}{N}
(G'WG)^{-1}G'WSWG(G'WG)^{-1}
\]

\textbf{Delta method}
\[
\widehat{\Var}(\psi(\widehat\theta))
\approx
\widehat D_\psi
\widehat{\Var}(\widehat\theta)
\widehat D_\psi'
\]

\textbf{Bootstrap}
\[
\widehat\theta^{*(b)}
=
\widehat\theta(Y_1^{*(b)},\ldots,Y_N^{*(b)})
\]
\end{frame}

\begin{frame}{Dependence and simulation error}
\begin{itemize}
    \item Cluster or block resampling should follow the independent source of variation.
    \item Panel histories need resampling schemes that preserve state paths.
    \item Market or policy moments often require market-level or policy-level clustering.
    \item SMM and indirect inference need seed checks and enough simulation draws.
    \item Generated regressors and two-step objects require correction or full-pipeline bootstrap.
    \item Report the uncertainty for the counterfactual object, not only primitive parameters.
\end{itemize}
\end{frame}

\begin{frame}{Theory to applied examples}
\scriptsize
\begin{tabular}{p{0.16\linewidth} p{0.24\linewidth} p{0.24\linewidth} p{0.24\linewidth}}
\toprule
Paper & Question & Estimation lesson & Fit / counterfactual lesson \\
\midrule
Rust & replacement timing & dynamic likelihood & hazards by state; cost counterfactuals \\
BLP / Nevo & differentiated demand & moments, instruments, simulation & substitution before policy claims \\
Todd-Wolpin & schooling subsidy & dynamic model with validation & experiment disciplines simulation \\
Low-Meghir-Pistaferri & lifecycle risk & moments for risk and insurance & lifecycle fit and welfare \\
Adda-Dustmann-Stevens & career costs of children & dynamic lifecycle structure & long-run state-specific paths \\
\bottomrule
\end{tabular}
\end{frame}

\begin{frame}{Limitations and credibility threats}
\begin{itemize}
    \item Identification by assumption: key primitives move because of normalization, not data.
    \item Weak moments: targeted objects do not move the primitive enough.
    \item Weighting dependence: conclusions change when moment weights change.
    \item Simulation noise: the optimizer chases random draws or local minima.
    \item Off-target failure: the model matches means but misses distributions or transitions.
    \item Counterfactual extrapolation: the policy changes support, menus, beliefs, or equilibrium conditions.
    \item Underreported uncertainty: first-stage, clustering, or simulation error is ignored.
\end{itemize}
\end{frame}

\begin{frame}{Research Lab design}
\begin{columns}[T,onlytextwidth]
\begin{column}{0.32\linewidth}
\textbf{Reproduce}
\begin{itemize}
    \item Rust-style replacement
    \item likelihood objective
    \item GMM moments
    \item targeted fit table
\end{itemize}
\end{column}
\begin{column}{0.32\linewidth}
\textbf{Diagnose}
\begin{itemize}
    \item weighting sensitivity
    \item untargeted support
    \item cluster bootstrap
    \item delta-method counterfactual
\end{itemize}
\end{column}
\begin{column}{0.32\linewidth}
\textbf{Transfer}
\begin{itemize}
    \item lifecycle schooling/work
    \item SMM moment fit
    \item tuition counterfactual
    \item validation memo
\end{itemize}
\end{column}
\end{columns}
\end{frame}

\begin{frame}{Takeaways}
\begin{itemize}
    \item Estimation strategy is part of the identification argument.
    \item Likelihood, moments, SMM, and indirect inference differ in what they require the model to explain.
    \item Weighting and targeting choices should be reported as substantive design choices.
    \item Fit is useful only when tied to the counterfactual and shown beyond construction.
    \item Inference must match the estimator, dependence structure, simulation design, and policy object.
\end{itemize}
\end{frame}

\end{document}
